\newpage
\begin{center}
  \textbf{\large 2. Анализ требований и постановка задачи разработки государственной информационной системы}
\end{center}
\refstepcounter{chapter}
\addcontentsline{toc}{chapter}{2. Анализ требований и постановка задачи разработки государственной информационной системы}

\section{Обзор существующих информационных систем управления цифровыми активами}

Существующие системы управления цифровыми активами можно разделить на несколько классов в зависимости от их функциональности и области применения. К основным классам относятся:

\begin{itemize}
  \item DAM-системы (Digital Asset Management) – предназначены для хранения, организации и управления цифровыми активами, такими как изображения, видео, аудио и документы. Примеры: Adobe Experience Manager, Bynder, Widen Collective.
  \item ECM / DMS-системы (Enterprise Content Management / Document Management Systems) – ориентированы на управление корпоративным контентом и документами. Примеры: Microsoft SharePoint, OpenText, Alfresco. Эти системы обеспечивают функции хранения, организации и управления документами, а также контроль версий и прав доступа.
  \item MAM-системы (Media Asset Management) – специализированные системы для управления медиа-активами, такими как видео и аудио. Примеры: Avid MediaCentral, CatDV, Dalet Galaxy.
  \item DRM / IAM (Digital Rights Management / Identity and Access Management) – системы, обеспечивающие контроль доступа и управление правами на цифровые активы.
\end{itemize}

Все вышеперечисленные системы разработаны для решения конкретных задач, в той или иной степени связанных с управлением цифровыми активами, но ни одна из них не предназначена специально для управления цифровым наследством. Они не обеспечивают специализированные функции, необходимые для управления цифровыми активами в рамках наследства, такие как автоматическое наследование прав доступа, интеграция с государственными реестрами или обеспечение юридической защиты участников процесса, а также невозможность доступа к данным до наступления триггерного события и долгосрочное хранение информации.

Более того, на текущий момент каждая из описанных выше систем, как правило, ориентирована на определённые типы цифровых активов и не обеспечивает универсального решения для управления всеми видами цифрового наследства.

Отметим, что при разработке системы частично могут быть использованы, например, функции проверки подлинности и контроля доступа, которые присутствуют в системах DRM / IAM, а также функции хранения и организации данных из DAM-систем.

Кроме того, существуют сервисные механизмы, встроенные в конкретные цифровые платформы и позволяющие пользователю заранее определить судьбу отдельных аккаунтов и данных. Юридический опыт, изученны в 1 главе работы показывает, что следует утвердить главенство воли наследодателя и, впервую очередь, использовать информацию из таких служб, тем более это позволит избегать проблем с нарушением условий использования сервиса внезависимости от его юрисдикции.

Отдельный интерес для предмета исследования представляет российская нотариальная инфраструктура. Единая информационная система нотариата позиционируется как централизованная и защищенная цифровая экосистема, обеспечивающая электронный документооборот, проверки сведений и доступ нотариусов к юридически значимой информации. Кроме того, нотариат уже поддерживает ряд удаленных действий, включая принятие на хранение и выдачу электронных документов, а также ведет реестр наследственных дел, позволяющий установить факт открытия наследственного дела и нотариуса, у которого оно находится. Это означает, что в российской правовой системе уже существует технологическая и институциональная основа для работы с электронными документами и наследственными процедурами, но пока отсутствует гражданско-ориентированный контур каталогизации и прижизненного управления цифровыми активами как объектом будущего наследования.

Также рассмотрим платформы, которые специально разработаны для управления цифровым наследством. Одним из примеров является платформа Inheriti®. Это коммерческая платформа, которая предоставляет пользователям возможность каталогизировать и передавать цифровые активы в рамках наследства, обеспечивая при этом безопасность и контроль доступа. Платформа работает следующим образом: пользователь описывает активы (логины для социальных сетей, электронные кошельки, файлы и т.д.) и указывает наследников, которые смогут получить доступ к этим активам после его смерти. Со своей стороны платформа обеспечивает безопасное хранение данных и передачу доступа наследникам.

Пользователь самостоятельно описывает свои цифровые активы и, что следует понимать, самостоятельно несёт ответственность не только за их актуальность и полноту, но и за любые юридические последствия, которые могут возникнуть в случае передачи доступа к этим активам наследникам. Это может включать в себя вопросы, связанные с авторским правом, конфиденциальностью и безопасностью данных. Например, передача аккаунта в социальной сети может привести к нарушению условий использования платформы, а передача доступа к электронному кошельку может вызвать проблемы с безопасностью и финансовыми потерями. Поэтому важно, чтобы пользователи были осведомлены о возможных рисках и юридических последствиях, связанных с управлением цифровым наследством, и принимали обоснованные решения при описании своих цифровых активов и выборе наследников.

Подобный сервис может частично решить задачу управления цифровым наследством, но он не является институциональным решением, поэтому не обеспечивает достаточный уровень унификации, стандартизации и не предполагает юридическую защиту для участников процесса.

